\documentclass[11pt,a4paper]{article}
\usepackage[utf8]{inputenc}
\usepackage[T1]{fontenc}
\usepackage{graphicx,booktabs,amsmath,hyperref}
\usepackage[margin=2.5cm]{geometry}
\graphicspath{{../results_extension/severity/}{../results_extension/localization/}{../results_extension/segmentation/}{../results_extension/}{figures/}}

% Tables come from `python extensions/make_paper_tables.py`. A missing table prints a visible note instead of failing.
\newcommand{\maybeinput}[1]{\IfFileExists{tables/#1}{\input{tables/#1}}{\textbf{[missing table: \detokenize{#1}]}}}

\title{TODO: working title}
\author{TODO: authors and affiliations}
\date{}

\begin{document}
\maketitle

\begin{abstract}
% Write this last, once the numbers are in. State the task, the data, the three additions, and the main
% result with its numbers, and say that severity labels are proxy labels and that localization and
% segmentation use weak labels.
\end{abstract}

\section{Introduction}
% Why automated crack screening matters. The gap in the original study: it classifies crack / no crack and
% explicitly leaves out severity, localization and segmentation. What this paper adds, in one paragraph.

\section{Related work}
% The original SSL study \cite{original_ssl_bridge}, DINOv2 \cite{dinov2}, U-Net \cite{unet}, YOLO
% \cite{ultralytics_yolo}, and published work on crack detection and segmentation that you have read.

\section{Data}
% SDNET2018 \cite{sdnet2018}: surfaces, image size, the number of cracked and uncracked images (take the
% counts from your own run). Photo-aware split: train / validation / test, why it is used, and how it
% differs from the random split of the original study.

\section{Methods}
\subsection{Stage 1: crack detection}
% Frozen DINOv2 ViT-S/14 features and a linear probe, as in the original study.

\subsection{Proxy severity labels}
% The measurement pipeline (severity/measure_cracks.py), the score, the cut-offs, how the cut-offs were
% fitted (training split only), and the spot-check against a human reviewer if you did one.

\subsection{Severity classifier}
% Frozen head with validation-selected regularisation, class-balanced loss, the optional fine-tuning of the
% last blocks, and the human-review threshold.

\subsection{Weak labels for localization and segmentation}
% How boxes and masks are generated from the measurement code, how negatives are added, and the preview figure.

\subsection{Localization and segmentation models}
% YOLO settings and the U-Net (architecture, loss, augmentation, threshold chosen on validation).

\section{Experiments}
% Metrics: per-class precision, recall, F1, macro-F1, confusion matrix, Severe-as-Minor count for severity;
% box precision, recall, mAP for localization; IoU and Dice for segmentation. Seeds and hardware.

\section{Results}
\subsection{Severity}
\maybeinput{severity_test.tex}
\maybeinput{severity_critical.tex}
\maybeinput{severity_frozen_vs_finetuned.tex}
% \includegraphics[width=0.5\linewidth]{severity_confusion_matrix.png}

\subsection{Localization}
\maybeinput{localization_test.tex}
% \includegraphics[width=\linewidth]{localization_examples.png}

\subsection{Segmentation}
\maybeinput{segmentation_test.tex}
% \includegraphics[width=0.7\linewidth]{segmentation_examples.png}

% Show failures as well as good cases, and explain them.

\section{Limitations}
% Image-only assessment cannot establish structural safety. Severity labels are proxies from measured crack
% extent. Boxes and masks are weak labels, so scores measure agreement with the measurement code. No
% physical scale calibration, so no real crack widths. Patch-level images from one dataset. Findings about
% generalisation to other bridges or cameras are not tested.

\section{Conclusion}
% What was done and what the numbers show, without going beyond them.

\bibliographystyle{plain}
\bibliography{references}
\end{document}
